\section{Discussion and Limitations}
\label{sec:discussion}

\textit{What each part of the loss does.} The study of the loss terms separates two roles. The
multi-view supervision carries most of the accuracy gain on the unseen ETIS set. The consistency term
and its label-aware weights complete the loss: with them, FCT is the only variant that improves all
five test sets on average, and the consistency gradient is twice as well aligned with supervision as
without the weights. By
Eq.~\eqref{eq:combined}, this division of roles is expected: under full supervision, the consistency
term brings no information beyond the label, only a pixel weight, so its effect on accuracy is that
of a reweighting of the supervised loss, and its weights decide whether that reweighting ever turns
against the label.

\textit{Why the unseen sets?} For all three networks, FCT gains most on ETIS, the test set from
another centre, and, for PraNet-V2, it also gains on CVC-ColonDB. These are the two sets on which every
network is least accurate, and on which Polyp-PVT and PraNet-V2 are least flip-consistent. One
explanation is the hypothesis of Section~\ref{sec:intro}: training on the flipped views and tying them
together discourages cues that change under a flip, and such cues are unlikely to transfer to a new
centre. Our data support this only indirectly. The flip gap falls on every set, but its fall does not
predict where Dice rises, and a single run with flip and rotation augmentation lowered it without any
gain.

\textit{Scope of the analysis.} Proposition~\ref{prop:identity} needs a binary label and a
cross-entropy or squared-error consistency term. It applies to each class of a multi-label task, one
class against the rest, but not directly to a softmax over several classes, where the two error vectors
need not be parallel and a consistency term can carry information beyond the label. Soft labels also
break it, since the two errors may then lie on opposite sides of the target; FCT therefore binarizes the
label at 0.5 when it computes the weights, including at rescaled input sizes, where interpolation
softens the mask boundary. The analysis also stops at the logits: it does not describe how the network
shares the per-pixel terms through its parameters.

\textit{Cost.} FCT multiplies the training time by about 3.5 and leaves inference unchanged. Without
the term on the upright view, the factor drops to three (Table~\ref{tab:cost}) at a similar accuracy
(Table~\ref{tab:rules}); applying FCT to part of each batch would lower it further.

\textit{Limits of the evidence.} The evidence has five limits. First, it covers three networks and one
training set of 1,288 images; of the five standard test sets, only ETIS comes from another centre, and
on the six PolypGen centres only the confidence-graded variant has been tested so far. Second, the
reported epoch is selected on the test sets. The selection is the same for both arms, and the gains
remain positive on average without it but are not significant at the last epoch
(Fig.~\ref{fig:curves}); the absolute scores are optimistic. Third, the weights were developed over
several preliminary variants scored on the same five test sets, so these sets also guided their design.
The confidence-graded weight, for instance, was carried to six seeds because of its result on seed~1,
which row (h) of Table~\ref{tab:rules} includes; on seeds 2--6 alone, its five-set mean gain is $+0.32$
$[0.08, 0.56]$. The gain of FCT is smallest on seed~1, $+0.18$ points, and is $+0.59$ $[0.15, 1.03]$ on
seeds 2--6. Fourth, six seeds resolve gains of about half a point on the five-set mean, but not the
smaller differences between variants of the loss, whose CIs are up to several points wide. Fifth,
$\lambda$ was set on Polyp-PVT and not retuned for the other networks, and the study of the loss terms
covers Polyp-PVT only, with one row still pending and two at three seeds.
